% GENERATED FILE. Edit canonical lessons, then run npm run generate:books.
% canonical-source-hash: fnv1a64:6dac2a7ab54e706e
% canonical-lessons: LA-C221-regula, LA-C221-discrimen, LA-C221-eventus, LA-C221-fortuna, LA-C221-casus

\chapter{Rule, Difference, Outcome, Fortune, Accident}
\label{ch:la-rule-difference-outcome-fortune-accident}
\hlchaptermodality{221}

\begin{chapteropening}
\textbf{By the end of this chapter:} \emph{I can say a rule, a difference, an outcome, fortune and an accident.}

Complete the last lesson of chapter 221: I can say a rule, a difference, an outcome, fortune and an accident.
\end{chapteropening}

\section[rēgula, rēgulae]{rēgula, rēgulae — a rule, a ruler}

\label{lesson:LA-C221-regula}



\begin{quote}
Before the new one: say the Latin for a young person, then the Latin for a widow.
\end{quote}

\subsection*{You'll want to know: rēgula, rēgulae}
\textbf{rēgula, rēgulae} — \textquotedblleft{}a rule, a ruler\textquotedblright{}.

\begin{cousinweb}
\textbf{rēgula} is built on \emph{regō}, \textquotedblleft{}to keep straight, to rule\textquotedblright{} (Lewis \& Short), with \emph{-ula}, an ending for a tool: a straight stick, a ruler. English \emph{rule} comes from it through French.
\end{cousinweb}

\begin{grammarlens}[title={What you've built}]
One more word for this chapter.
\end{grammarlens}

\subsection*{Your turn}
\begin{itemize}
\raggedright
  \item \emph{Say it:} \emph{rēgula, rēgulae}
  \item \emph{Say it:} \emph{rēgula, rēgulae}, once more
  \item \emph{Say it:} say \emph{rēgula, rēgulae}
  \item \emph{Recall:} say the Latin for a young person, then the Latin for a widow, then say \emph{rēgula, rēgulae} again
\end{itemize}

\begin{tcolorbox}[breakable,colback=teal!4,colframe=teal!35!black,title={Before you move on}]
What does rēgula, rēgulae mean? (\textquotedblleft{}A rule, a ruler\textquotedblright{}.) Say it once more.
\end{tcolorbox}
\section[discrīmen, discrīminis]{discrīmen, discrīminis — a difference; a crisis}

\label{lesson:LA-C221-discrimen}



\begin{quote}
Before the new one: say the Latin for a widow, then the Latin for a rule.
\end{quote}

\subsection*{You'll want to know: discrīmen, discrīminis}
\textbf{discrīmen, discrīminis} — \textquotedblleft{}a difference; a crisis\textquotedblright{}.

\begin{cousinweb}
\textbf{discrīmen} is contracted from older \emph{discerimen} (Lewis \& Short), from \emph{discernō}, \textquotedblleft{}to separate, to distinguish\textquotedblright{}, with \emph{-men}, an ending for a result: a dividing line, then a turning point.
\end{cousinweb}

\begin{grammarlens}[title={What you've built}]
One more word for this chapter.
\end{grammarlens}

\subsection*{Your turn}
\begin{itemize}
\raggedright
  \item \emph{Say it:} \emph{discrīmen, discrīminis}
  \item \emph{Say it:} \emph{discrīmen, discrīminis}, once more
  \item \emph{Say it:} say \emph{discrīmen, discrīminis}
  \item \emph{Recall:} say the Latin for a widow, then the Latin for a rule, then say \emph{discrīmen, discrīminis} again
\end{itemize}

\begin{tcolorbox}[breakable,colback=teal!4,colframe=teal!35!black,title={Before you move on}]
What does discrīmen, discrīminis mean? (\textquotedblleft{}A difference; a crisis\textquotedblright{}.) Say it once more.
\end{tcolorbox}
\section[ēventus, ēventūs]{ēventus, ēventūs — an outcome}

\label{lesson:LA-C221-eventus}



\begin{quote}
Before the new one: say the Latin for a rule, then the Latin for a difference; a crisis.
\end{quote}

\subsection*{You'll want to know: ēventus, ēventūs}
\textbf{ēventus, ēventūs} — \textquotedblleft{}an outcome\textquotedblright{}.

\begin{cousinweb}
\textbf{ēventus} is built on \emph{ēveniō}, \textquotedblleft{}to come out, to happen\textquotedblright{} (Lewis \& Short), from \emph{ē}, \textquotedblleft{}out\textquotedblright{}, and \emph{veniō}, \textquotedblleft{}to come\textquotedblright{}: how things come out.
\end{cousinweb}

\begin{grammarlens}[title={What you've built}]
One more word for this chapter.
\end{grammarlens}

\subsection*{Your turn}
\begin{itemize}
\raggedright
  \item \emph{Say it:} \emph{ēventus, ēventūs}
  \item \emph{Say it:} \emph{ēventus, ēventūs}, once more
  \item \emph{Say it:} say \emph{ēventus, ēventūs}
  \item \emph{Recall:} say the Latin for a rule, then the Latin for a difference; a crisis, then say \emph{ēventus, ēventūs} again
\end{itemize}

\begin{tcolorbox}[breakable,colback=teal!4,colframe=teal!35!black,title={Before you move on}]
What does ēventus, ēventūs mean? (\textquotedblleft{}An outcome\textquotedblright{}.) Say it once more.
\end{tcolorbox}
\section[fortūna, fortūnae]{fortūna, fortūnae — fortune, luck}

\label{lesson:LA-C221-fortuna}



\begin{quote}
Before the new one: say the Latin for a difference; a crisis, then the Latin for an outcome.
\end{quote}

\subsection*{You'll want to know: fortūna, fortūnae}
\textbf{fortūna, fortūnae} — \textquotedblleft{}fortune, luck\textquotedblright{}.

\begin{cousinweb}
\textbf{fortūna} is built on \emph{fors}, \textquotedblleft{}chance, luck\textquotedblright{} (Lewis \& Short), with \emph{-tūna}, an ending also in \emph{Neptūnus}, \textquotedblleft{}Neptune\textquotedblright{}. \emph{fors} is usually traced to the root of \emph{ferō}, \textquotedblleft{}to bring\textquotedblright{}: what chance brings.
\end{cousinweb}

\begin{grammarlens}[title={What you've built}]
One more word for this chapter.
\end{grammarlens}

\subsection*{Your turn}
\begin{itemize}
\raggedright
  \item \emph{Say it:} \emph{fortūna, fortūnae}
  \item \emph{Say it:} \emph{fortūna, fortūnae}, once more
  \item \emph{Say it:} say \emph{fortūna, fortūnae}
  \item \emph{Recall:} say the Latin for a difference; a crisis, then the Latin for an outcome, then say \emph{fortūna, fortūnae} again
\end{itemize}

\begin{tcolorbox}[breakable,colback=teal!4,colframe=teal!35!black,title={Before you move on}]
What does fortūna, fortūnae mean? (\textquotedblleft{}Fortune, luck\textquotedblright{}.) Say it once more.
\end{tcolorbox}
\section[cāsus, cāsūs]{cāsus, cāsūs — an accident, a fall}

\label{lesson:LA-C221-casus}



\begin{quote}
Before the new one: say the Latin for an outcome, then the Latin for fortune.
\end{quote}

\subsection*{You'll want to know: cāsus, cāsūs}
\textbf{cāsus, cāsūs} — \textquotedblleft{}an accident, a fall\textquotedblright{}.

\begin{cousinweb}
\textbf{cāsus} is built on \emph{cadō}, \textquotedblleft{}to fall\textquotedblright{} (Lewis \& Short): a falling, then what befalls. English \emph{case} comes from it.
\end{cousinweb}

\begin{grammarlens}[title={What you've built}]
One more word for this chapter.
\end{grammarlens}

\subsection*{Your turn}
\begin{itemize}
\raggedright
  \item \emph{Say it:} \emph{cāsus, cāsūs}
  \item \emph{Say it:} \emph{cāsus, cāsūs}, once more
  \item \emph{Say it:} say \emph{cāsus, cāsūs}
  \item \emph{Recall:} say the Latin for an outcome, then the Latin for fortune, then say \emph{cāsus, cāsūs} again
\end{itemize}

\begin{tcolorbox}[breakable,colback=teal!4,colframe=teal!35!black,title={Before you move on}]
What does cāsus, cāsūs mean? (\textquotedblleft{}An accident, a fall\textquotedblright{}.) Say it once more.
\end{tcolorbox}
